\section{Вычислительная эффективность}
\label{sec:performance}

\draft{Измерения: Intel Core i7-12700H, один вычислительный поток,
GCC 15.2, оптимизированная сборка, PETSc 3.24.4; $\Delta x=200$ м, $\Delta z=50$ м,
$\Delta t=1$ млн лет, $\varepsilon=0{,}05$ м, допуск упрощения 5 м.
Рисунок~\ref{fig:timing} показывает распределение времени по этапам.
Подготовка входа и экспорт XLSX в Python в эти времена не входят.}
\reportfigure{21_forward_stage_timings.png}
  {Доли времени прямого расчёта на трёх наборах:
  стационарный старт, сетка, перенос, сборка, линейное решение, Ro и вывод.}
  {fig:timing}

\draft{На рис.~\ref{fig:solvers} GMRES(100)/ILU(2) сопоставлен с прямым
решателем UMFPACK~\cite{suitesparse}. По медианам трёх запусков ускорение
линейного решения составляет 4,05; 3,67 и 1,44 раза для data1, data2 и обратного
набора соответственно; всего прямого расчёта --- 2,55; 2,32 и 1,31 раза.
Различия T между решателями не превышают $3\cdot10^{-11}\,{}^\circ$C.}
\reportfigure{22_linear_solver_comparison.png}
  {Время линейного решения с подготовкой и полного прямого расчёта:
  GMRES(100)/ILU(2) и UMFPACK, медианы и диапазоны трёх запусков.}
  {fig:solvers}

\draft{На рис.~\ref{fig:easy-speed} изолированное вычисление интегралов
EASY\%Ro с таблицей быстрее прямой аналитики через интегральную экспоненту
в 6,52 раза. Однократное построение таблицы занимает около 0,25 мс.
Это ускорение вычислительного ядра, а не всего прямого расчёта.}
\reportfigure{23_easy_ro_table_benchmark.png}
  {Стоимость вычисления интегралов через аналитическую первообразную
  и таблицу; справа --- максимальная относительная ошибка Ro
  на четырёх историях относительно независимой квадратуры.}
  {fig:easy-speed}

\draft{Рисунок~\ref{fig:rdp} показывает ускорение расчёта за счёт
упрощения геометрии алгоритмом RDP и связанное с ним изменение решения на data2:
при допуске 5 м число ячеек уменьшается с 120335 до 3943,
время --- с 205,4 до 5,0 с. Относительные табличные отличия средних T/Ro
равны 1,31/1,12\%, но локальное отличие Ro достигает 6,91\%.
Это одиночные замеры и сравнение с подробным численным расчётом;
геометрия и сетка меняются одновременно.}
\reportfigure{24_geometry_simplification_tradeoff.png}
  {Data2: влияние допуска упрощения границ на размер сетки,
  время и отличия современных T и Ro от расчёта без упрощения.}
  {fig:rdp}
